%% ch10_infrastructure.tex -- Intelligence Engineering, chapter 10
%% "Infrastructure and Tooling for MLOps".
%% Spine transferred from Huyen, Designing Machine Learning Systems, chapter
%% 10. The subchapters and the frame titles under them are the book's own
%% section and subsection headings, in the book's order.
%%
%% STATE: titles and TEMPLATE layout only. No body text. Every frame carries
%% the layout it is meant to be filled into, and a % TEMPLATE comment saying
%% why that layout and what belongs in each slot. Filling them is the next pass.
%%
%% Fragment of intelligence_engineering.tex. One chapter is one lecture and
%% gets exactly ONE \lecturepage, at the top. Inside it every \subsection
%% opens on the subchapter divider, which the master emits automatically via
%% \AtBeginSubsection; do not write those divider frames by hand.

\begin{refsection}

\lecturepage{X}{Infrastructure and Tooling for MLOps}{}
\section[Infrastructure]{Infrastructure and Tooling for MLOps}

% ----------------------------------------------------------------------
\subsection{Storage and Compute}

% TEMPLATE: two block contrast, levelled by the equal height group. The heading
% is a versus, so each side takes one box rather than a bullet list and the box
% headers are read straight off the two sides the heading names. Public cloud
% box takes elasticity and the start without capital outlay; private box takes
% what owning the machines buys back, with the book's cloud bill numbers, the
% revenue share and the repatriation saving, kept as numbers.
\begin{frame}{Public Cloud Versus Private Data Centers}
  \begin{columns}[T,onlytextwidth]
    \begin{column}{0.47\textwidth}
      \begin{tdblockt}[equal height group=cloud]{Public cloud}
      \end{tdblockt}
    \end{column}
    \begin{column}{0.47\textwidth}
      \begin{tdblockt}[equal height group=cloud]{Private data centers}
      \end{tdblockt}
    \end{column}
  \end{columns}
  \slidesources{Huyen2022}
\end{frame}

% ----------------------------------------------------------------------
\subsection{Development Environment}

% TEMPLATE: captioned code panel, kept bare. The book's dev environment is a
% list of tools and the configuration that pins them, the IDE, versioning and
% CI/CD, so the honest slot is a panel and not prose about tooling. Caption is
% the heading's own subject; the panel takes the named tools and the versioning
% setup in the book's order.
\begin{frame}{Dev Environment Setup}
  \begin{tdcodet}{Dev environment}
  \end{tdcodet}
  \slidesources{Huyen2022}
\end{frame}

% TEMPLATE: withmargin, the chapter's workhorse and the right shape here: one
% claim plus one aside. Main column takes the book's recommendation, standardize
% the environment company wide and move development to the cloud, and what that
% move costs. Margin takes the resistance the book records and the remote IDE it
% names as the way around it.
\begin{frame}{Standardizing Dev Environments}
  \begin{withmargin}
    \begin{tdmain}
    \end{tdmain}
    \begin{tdmargin}
    \end{tdmargin}
  \end{withmargin}
  \slidesources{Huyen2022}
\end{frame}

% TEMPLATE: withmargin with a captioned code panel in the main column. The
% heading is about containers, which are configuration before they are an idea,
% so the panel takes the Dockerfile, the image against container distinction and
% the orchestration tools in the book's words. Margin takes why the same
% environment has to travel from dev to prod at all.
\begin{frame}{From Dev to Prod: Containers}
  \begin{withmargin}
    \begin{tdmain}
      \begin{tdcodet}{Containers}
      \end{tdcodet}
    \end{tdmain}
    \begin{tdmargin}
    \end{tdmargin}
  \end{withmargin}
  \slidesources{Huyen2022}
\end{frame}

% ----------------------------------------------------------------------
\subsection{Resource Management}

% TEMPLATE: three block row, headers taken from the three peers the heading
% names, levelled by the equal height group. The book's point is that these
% three are not the same thing, which a row of boxes shows and a bullet list
% hides. Each box takes what that tool schedules and what it cannot do, with
% the book's division of labour, jobs against resources, landing in the last.
\begin{frame}{Cron, Schedulers, and Orchestrators}
  \begin{columns}[T,onlytextwidth]
    \begin{column}{0.31\textwidth}
      \begin{tdblockt}[equal height group=rm]{Cron}
      \end{tdblockt}
    \end{column}
    \begin{column}{0.31\textwidth}
      \begin{tdblockt}[equal height group=rm]{Schedulers}
      \end{tdblockt}
    \end{column}
    \begin{column}{0.31\textwidth}
      \begin{tdblockt}[equal height group=rm]{Orchestrators}
      \end{tdblockt}
    \end{column}
  \end{columns}
  \slidesources{Huyen2022}
\end{frame}

% TEMPLATE: side image hero, the chapter's one and only figure. A workflow is a
% graph of steps and the book draws it as one, so the strip takes that figure
% and the sidebody takes the book's comparison of the named workflow managers,
% what each of them defines a step in and what it does about containers. The
% caption names the figure; the gray placeholder stays until the asset is drawn.
\begin{frame}{Data Science Workflow Management}
  \phimgside[0.33]{data_science_workflow_management.png}{Data science workflow}
  \begin{sidebody}
  \end{sidebody}
  \slidesources{Huyen2022}
\end{frame}

% ----------------------------------------------------------------------
\subsection{ML Platform}

% TEMPLATE: withmargin. The heading is one platform component and not a set of
% peers, so the main column runs it as a single claim: what the deployment
% component owes the rest of the platform once a model is pushed, online and
% batch prediction behind one service. Margin takes the named deployment tools
% the book lists and the caveat it attaches to buying them.
\begin{frame}{Model Deployment}
  \begin{withmargin}
    \begin{tdmain}
    \end{tdmain}
    \begin{tdmargin}
    \end{tdmargin}
  \end{withmargin}
  \slidesources{Huyen2022}
\end{frame}

% TEMPLATE: tddef, kept bare. Model store is a term the book has to define
% before it can show that most teams implement only a fraction of it, so the
% term box is the whole slide rather than a claim with an aside. The box takes
% the definition and then the artifacts the book says belong in a model store,
% which is the part the tools skip.
\begin{frame}{Model Store}
  \begin{tddef}{Model store}
  \end{tddef}
  \slidesources{Huyen2022}
\end{frame}

% TEMPLATE: three block row, untitled on purpose. The book settles this heading
% by naming three problems a feature store solves, so three levelled boxes
% rather than three bullets in one; the headers stay off because the heading
% names only the store itself and the three names arrive with the text in the
% next pass. Each box takes one problem and the tools the book credits with it.
\begin{frame}{Feature Store}
  \begin{columns}[T,onlytextwidth]
    \begin{column}{0.31\textwidth}
      \begin{tdblock}[equal height group=fs]
      \end{tdblock}
    \end{column}
    \begin{column}{0.31\textwidth}
      \begin{tdblock}[equal height group=fs]
      \end{tdblock}
    \end{column}
    \begin{column}{0.31\textwidth}
      \begin{tdblock}[equal height group=fs]
      \end{tdblock}
    \end{column}
  \end{columns}
  \slidesources{Huyen2022}
\end{frame}

% ----------------------------------------------------------------------
\subsection{Build Versus Buy}

% TEMPLATE: booktabs comparison table. This is the chapter's second versus and
% the two box contrast is already spent on the cloud frame, so the decision goes
% into a table: two columns headed by the two sides the heading names, the book's
% deciding questions as the rows. Row labels and cells are filled together in
% the next pass, which is why both are empty here.
\begin{frame}{Build Versus Buy}
  \footnotesize
  \renewcommand{\arraystretch}{1.35}
  \begin{tabular}{@{}%
    >{\raggedright\arraybackslash}p{0.22\textwidth}%
    >{\raggedright\arraybackslash}p{0.35\textwidth}%
    >{\raggedright\arraybackslash}p{0.35\textwidth}@{}}
    \toprule
     & \textbf{Build} & \textbf{Buy}\\
    \midrule
     & & \\
     & & \\
     & & \\
    \bottomrule
  \end{tabular}
  \slidesources{Huyen2022}
\end{frame}

% ----------------------------------------------------------------------
\subsection{Summary}

% TEMPLATE: takeaway, kept bare. A summary is a closing line and not a slide of
% bullets, so it takes the big centred form and nothing else; the chapter has no
% other takeaway, so this one lands. The line is the book's own closing claim,
% that what a team needs from infrastructure is set by the number and the scale
% of its use cases.
\begin{frame}{Summary}
  \begin{takeaway}
  \end{takeaway}
  \slidesources{Huyen2022}
\end{frame}

% ----------------------------------------------------------------------
% This chapter's own references. The refsection opened above scopes the
% citation numbering to this chapter, so chapter_pdfs/<this>.pdf resolves
% its own [n] markers instead of pointing at a list it does not contain.
\begin{frame}{References}
  \printbibliography[heading=none]
\end{frame}
\end{refsection}
